\chapter{Appendices}\label{sec:appendices}
\section{Appendix B: Model Architectures and Hyperparameters}\label{sec:appendix-b-model-architectures-and-hyperparameters}

% column helpers
\newcolumntype{L}[1]{>{\raggedright\arraybackslash}p{#1}}
\newcolumntype{Y}{>{\raggedright\arraybackslash}X}

\begin{table}[ht]
\centering
\footnotesize
\setlength{\tabcolsep}{4pt}
\renewcommand{\arraystretch}{1.15}

\begin{tabularx}{\textwidth}{|L{2.1cm}|Y|L{3.2cm}|L{1.7cm}|L{2.3cm}|}
\hline
\textbf{Relay} &
\textbf{Network Topology} &
\textbf{Train Loss (LR)} &
\textbf{Total Params} &
\textbf{Implemen-}\\
\textbf{tation} \\
\hline
MLP (Minimal) &
Feedforward MLP, 5-symbol sliding window &
MSE, lr=0.01, 100 epochs &
169 &
NumPy (CPU) \\
\hline
Hybrid &
MLP + DF switching via learned SNR threshold ($\sim$5 dB) &
MSE (MLP only), lr=0.01 &
169 &
NumPy (CPU) \\
\hline
VAE &
Encoder--decoder latent model (7$\to$8$\to$1) &
$\beta$-VAE ($\beta=0.1$), Adam lr=$10^{-3}$ &
1{,}777 &
PyTorch (CUDA) \\
\hline
cGAN (WGAN-GP) &
Conditional GAN (generator + critic) &
WGAN-GP, Adam lr=$10^{-4}$ &
2{,}946 &
PyTorch (CUDA) \\
\hline
Transformer &
2-layer encoder, 4 heads, $d_{\text{model}}=32$ &
MSE, Adam lr=$10^{-3}$ &
17{,}697 &
PyTorch (CUDA/CPU) \\
\hline
Mamba-S6 &
Selective SSM, 2 S6 layers &
MSE, Adam lr=$10^{-3}$ &
24{,}001 &
PyTorch (CUDA/CPU) \\
\hline
Mamba-2 SSD &
Selective SSM with SSD kernel (2 layers) &
MSE, Adam lr=$10^{-3}$, clip=1.0 &
26{,}179 &
PyTorch (CUDA) \\
\hline
\end{tabularx}

\caption[Relay models, objectives and parameter counts]{Comparison of relay models, training objectives, parameter counts, and implementations}
\label{tab:relay_comparison}
\end{table}

\section{Appendix C: Reproducibility}\label{sec:appendix-c-software-architecture}

\subsection{Testing} The simulation framework carries 187 automated pytest tests covering channels, modulation, relay strategies, simulation, and statistics. In the recorded environment, the full suite passed; with PyTorch unavailable, 183 passed and four checkpoint-dependent tests were skipped.

\subsection{Reproducibility} Evaluation seeds are controlled for source-bit and noise generation. Training initialization was not consistently seeded for every original sequence-model checkpoint, so exact retraining of those models is not guaranteed; the later replicated sweeps report training-seed variation explicitly.

\paragraph{Provenance of the reported results.} Each reported result is associated with an experiment artifact and a verification script in the source-code repository, using the experiment-specific budgets documented in the relevant section and ledger. Independent audit requires that repository at the archived revision identifier; the ledger below documents the claimed provenance but cannot by itself reproduce every table cell.

\clearpage

\section{Appendix D: Master Experiment Ledger}\label{sec:appendix-e-master-experiment-ledger}

This ledger links major studies to their experiment artifacts. The table verifier checks a declared subset of numerical cells, not every number or figure. Historical \texttt{e6\_*.npy} paths are relative to \url{e6_unknown_channel_results/}; the central ISI figure and table additionally use the matched-protocol and fixed-budget companion files listed below.

{\footnotesize
\begin{longtable}[]{@{}p{0.18\linewidth}p{0.10\linewidth}p{0.12\linewidth}p{0.12\linewidth}p{0.36\linewidth}@{}}
\caption[Master experiment ledger]{Master experiment ledger: modulation, channel, and result-file provenance for every major study}\label{tbl:ledger}\tabularnewline
\toprule\noalign{}
Study & \shortstack[l]{Modu-\\lation} & Channel & Tbl/Fig & Result source \\
\midrule\noalign{}
\endfirsthead
\toprule\noalign{}
Study & \shortstack[l]{Modu-\\lation} & Channel & Tbl/Fig & Result source \\
\midrule\noalign{}
\endhead
\bottomrule\noalign{}
\endlastfoot
Canonical benchmark & QPSK & Rayleigh & tbl:2, fig:10 & \url{results/modulation/qpsk_rayleigh.json} \\
Normalized 3K-param.\ comparison & QPSK & Rayleigh & tbl:8 & \url{results/normalized_3k/3k_rayleigh.json} \\
Coded block-DF, $K$-sweep & QPSK & Rayleigh & tbl:34, 35 & \url{results/coded_df_experiment.json} \\
Coded high-SNR paired re-test & QPSK & Rayleigh & tbl:37 & \url{results/coded_high_budget_test.json} \\
Coded relay-error mechanism & QPSK & Rayleigh & tbl:38 & \url{results/coded_error_mechanism.json} \\
Coded soft-decision relaying & QPSK & Rayleigh & tbl:39 & \url{results/coded_soft_decision.json} \\
Equal-throughput \& latency & QPSK, 16-QAM & Rayleigh & tbl:40, 41 & \url{results/coded_latency_throughput.json} \\
Adaptive modulation/coding (AMC) & QPSK, 16-QAM & Rayleigh & tbl:42 & \url{results/coded_rate_adaptation.json} \\
Capacity under a latency budget & QPSK, 16-QAM & Rayleigh & tbl:43 & \url{results/coded_latency_capacity.json} \\
Unknown ISI (layer 2) & BPSK & AWGN (Hop 2) & tbl:E6 & Historical AF/DF/MLP below 12~dB: \url{e6_sim_ported_results.npy}; fixed-budget companion \url{codex_isi_fixed_budget_validation.json}; Viterbi: \url{results/e6_matched_protocol.json} and \url{results/e6_matched_highsnr.json}. \\
Flat-channel control & BPSK & Rayleigh & tbl:E6flat & \url{e6_flat_ported_results.npy} \\
Unknown ISI, QPSK generalization & QPSK & AWGN & \shortstack[l]{E6\\QPSK} & \url{e6_qpsk_unknown_channel_results.npy} \\
Composite impairment cascade & DBPSK & Rayleigh & \shortstack[l]{E6\\composite} & \url{e6_composite_ported_results.npy}; AF replaced by \url{codex_composite_af_validation.json}. \\
Partial posterior (pilot sweep) & BPSK & Rayleigh & E6partial & \url{e6_partial_ported_results.npy} \\
Blind channel (layer 4) & BPSK & Rayleigh & E6blind & \url{e6_blind_ported_results.npy} \\
Four-layer summary & mixed (above) & mixed (above) & tbl:layers & combination of the above \\
\end{longtable}
}

The result files above are themselves generated by named, committed scripts (e.g.\ \url{coded_rate_adaptation.py} produces \url{coded_rate_adaptation.json}); the correspondence between script and output file is one-to-one and is not separately tabulated here since the filenames make it explicit.

\begin{center}\rule{0.5\linewidth}{0.5pt}\end{center}

%% ── Hebrew Abstract (Back Matter) ───────────────────────────

\section{Appendix E: Minimum Relay Size and Initialization Variance}\label{sec:appendix-f-minimum-size}

Supporting material for the minimum-size study of Section~\ref{sec:minimum-relay-size} and the initialization-variance measurements discussed in Section~\ref{sec:capacity-and-overfitting}. The findings themselves are stated in those sections; what follows is the underlying tabulation and the figures, collected here to keep the experiment and discussion chapters to their conclusions.

\begin{longtable}[]{@{}llllllll@{}}
\caption[Smallest relay per degradation budget]{Smallest relay meeting a given degradation budget against each scenario's own published relay (MLP-169, or MLP-756 for the coded row), written as parameters and window. Most memoryless rows settle at window~1, but flat gain and nonlinear bias require wider windows at the strictest budget; every evaluated three-tap row requires window~5 or~7 until a full decibel is conceded. A dash indicates that no smaller configuration meets that budget}\label{tbl:table-minsize}\tabularnewline
\toprule\noalign{}
Channel & Taps & $\le$0~dB & $\le$0.1~dB & $\le$0.25~dB & $\le$0.5~dB & $\le$1~dB & $\le$2~dB \\
\midrule\noalign{}
\endfirsthead
\toprule\noalign{}
Channel & Taps & $\le$0~dB & $\le$0.1~dB & $\le$0.25~dB & $\le$0.5~dB & $\le$1~dB & $\le$2~dB \\
\midrule\noalign{}
\endhead
\bottomrule\noalign{}
\endlastfoot
awgn & 1 & 4p w1 & 4p w1 & 4p w1 & 4p w1 & 4p w1 & 4p w1 \\
rayleigh & 1 & 4p w1 & 4p w1 & 4p w1 & 4p w1 & 4p w1 & 4p w1 \\
flat gain & 1 & 21p w3 & 21p w3 & 21p w3 & 4p w1 & 4p w1 & 4p w1 \\
branch asym. & 1 & 4p w1 & 4p w1 & 4p w1 & 4p w1 & 4p w1 & 4p w1 \\
nonlinear bias & 1 & 29p w5 & 4p w1 & 4p w1 & 4p w1 & 4p w1 & 4p w1 \\
ISI & 3 & 73p w7 & 73p w7 & 57p w5 & 57p w5 & 37p w7 & 21p w3 \\
ISI (complex) & 3 & 73p w7 & 73p w7 & 73p w7 & 73p w7 & 57p w5 & 21p w3 \\
ISI + Rayleigh & 3 & 145p w7 & 73p w7 & 73p w7 & 73p w7 & 37p w7 & 21p w3 \\
composite & 3 & 145p w7 & 73p w7 & 57p w5 & 21p w3 & 21p w3 & 6p w3 \\
coded & code & 372p w9 & 244p w5 & 244p w5 & 18p w1 & 18p w1 & 18p w1 \\
\end{longtable}

\subsection{Results Figures}\label{sec:results-figures-minsize}

\begin{figure}[htbp]
\centering
\pandocbounded{\includegraphics[width=\linewidth,height=0.8\textheight,keepaspectratio]{results/minsize_window_crossover.png}}
\caption[Window against channel memory]{SNR penalty against each channel's classical comparator as the relay window widens, for the best configuration at each width. The reader should observe that the two panels behave oppositely: on the memoryless channels the penalty is flat to within a fifth of a decibel across every window (see inset), so the window buys nothing, while on the three-tap channels it falls by 8 to 11~dB between window~1 and window~7. This is the sense in which the relay's window must be matched to the channel's impulse response}\label{fig:minsize-crossover}
\end{figure}

\begin{figure}[htbp]
\centering
\pandocbounded{\includegraphics[width=\linewidth,height=0.8\textheight,keepaspectratio]{results/minsize_vs_degradation_budget.png}}
\caption[Smallest relay per degradation budget]{Smallest relay meeting a given degradation budget against MLP-169, drawn as a step function because the horizontal axis is a budget the designer chooses rather than a measured quantity. Most memoryless cases reach four parameters once a small degradation is allowed, while strict zero-degradation budgets require larger models for flat gain and nonlinear bias; the three-tap channels need between 21 and 145 depending on the permitted degradation}\label{fig:minsize-budget}
\end{figure}

\begin{table}[htbp]
\centering
\caption[Across-initialization spread]{Spread in SNR penalty between the best and worst of three initializations at an identical configuration, summarized over all configurations swept on each channel. The reader should observe that on five of the nine channels the median spread is comparable to or larger than the entire effect of varying model size, so a single-initialization size comparison on those channels measures the draw rather than the architecture.}\label{tbl:seed-spread}
\begin{tabular}{lrrrr}
\toprule
Channel & Configs & Median (dB) & 90th pct.\ (dB) & Max (dB) \\
\midrule
rayleigh      & 17 & 0.008 & 0.019 & 0.028 \\
composite     & 16 & 0.036 & 0.165 & 0.195 \\
awgn          & 17 & 0.055 & 0.116 & 0.138 \\
isi\_rayleigh & 16 & 0.177 & 0.588 & 0.795 \\
isi           & 18 & 0.195 & 0.705 & 0.996 \\
isi\_complex  & 17 & 0.332 & 2.325 & 3.182 \\
branch\_asym  & 17 & 0.388 & 0.892 & 1.013 \\
flat\_gain    & 17 & 0.622 & 1.000 & 1.250 \\
nlbias        & 17 & 1.231 & 1.242 & 1.245 \\
\midrule
All           & 152 & 0.180 & 1.229 & 3.182 \\
\bottomrule
\end{tabular}
\end{table}

\begin{table}[htbp]
\centering
\caption[Across-initialization spread at the 3K budget]{Spread in SNR penalty against symbol-wise DF between the best and worst of three initializations, for the equal-budget architectures on the canonical channel. The reader should observe the Transformer row: its spread is five times the roughly $0.18$~dB that H4's $4.1\%$ high-SNR architecture spread corresponds to, so a single run of that architecture cannot resolve the effect H4 claims. Every other architecture is comfortably inside it}\label{tbl:seed-spread-3k}
\begin{tabular}{lrrr}
\toprule
Architecture & Params & Per-seed penalty vs DF (dB) & Spread (dB) \\
\midrule
Mamba2-3K       & 3{,}004 & $+1.113$, $+1.094$, $+1.098$ & 0.019 \\
Mamba-S6-3K     & 3{,}027 & $+1.123$, $+1.110$, $+1.082$ & 0.041 \\
MLP-3K          & 3{,}004 & $+0.133$, $+0.162$, $+0.189$ & 0.055 \\
VAE-3K          & 3{,}037 & $+0.072$, $+0.228$, $+0.142$ & 0.156 \\
Transformer-3K  & 3{,}007 & $+2.039$, $+1.147$, $+1.133$ & \textbf{0.906} \\
\bottomrule
\end{tabular}
\end{table}

Eight Transformer initializations are shown in the following figures: seven cluster and one has both higher final loss and worse high-SNR BER. A sample of three cannot reliably characterize such failures. If three runs were sampled without replacement from these eight, the probability of missing the one outlier would be $\binom{7}{3}/\binom{8}{3}=5/8$; this finite-set calculation is not an estimate of the population failure rate. The association does not validate training loss as a general screening rule.

\begin{figure}[htbp]
\centering
\pandocbounded{\includegraphics[width=0.78\linewidth,height=0.34\textheight,keepaspectratio]{results/transformer_seed_ber_curves.png}}
\caption[Transformer-3K across eight initializations]{Transformer-3K on the canonical channel across eight initializations, against symbol-wise DF. The reader should observe that the failing run is not offset by a constant: it tracks the others at $0$~dB and separates as SNR rises, so an evaluation confined to low SNR would not distinguish it from the seven that converge}\label{fig:transformer-seed-curves}
\end{figure}

\begin{figure}[htbp]
\centering
\pandocbounded{\includegraphics[width=0.78\linewidth,height=0.34\textheight,keepaspectratio]{results/transformer_loss_vs_penalty.png}}
\caption[Training loss against SNR penalty]{Final training loss against SNR penalty for the eight Transformer initializations and the four-seed Mamba-S6 control. The correlation over all eight seeds is $r=+0.998$, but that figure is the outlier's leverage: among the seven runs that converge it is $r=-0.38$ ($p=0.41$), so loss identifies the failure and carries no information about relative quality once training has succeeded}\label{fig:transformer-loss-penalty}
\end{figure}

\section{Appendix F: Verification of Headline Claims}\label{sec:appendix-claim-verification}

Each headline claim of Chapter~\ref{sec:unknown-channels} is restated here with the argument or derivation that supports it and the primary reference against which it was checked. The channel used throughout is the normalized 3-tap FIR $\mathbf{h} = [1, 0.7, 0.5]/\lVert\cdot\rVert \approx [0.758, 0.531, 0.379]$.

\subsection{Matched-receiver bounds and the empirical MLSE comparison}

For a fixed observation and supplied model, bit-MAP minimizes expected bit error; MLSE instead minimizes sequence error \cite{Forney1972MLSE,BahlCockeJelinekRaviv1974BCJR,ProakisSalehi2008DigitalComm}. Thus a lower measured BER than MLSE would not contradict its optimality. The QPSK BCJR control measures relay-output errors only and must not be compared directly with destination BER. It does not close an end-to-end BCJR benchmark. Learned-receiver studies use different information assumptions and objectives \cite{FarsadGoldsmith2018NNDetection,Shlezinger2020ViterbiNet,Honkala2021DeepRx,YeLiJuang2018OFDMDL,OSheaHoydis2017PhysicalLayer}; they should not be summarized as a universal claim that learning wins only under mismatch.

\subsection{The analytic 0.25 floor and the non-monotonicity of symbol-wise DF}

With the normalized taps, the interferer magnitude sum $h_1 + h_2 \approx 0.910$ exceeds the cursor $h_0 \approx 0.758$: the eye is closed. Of the four equiprobable BPSK interferer sign patterns, exactly one (both opposing the cursor) yields a noise-free sample of $0.758 - 0.910 = -0.152 < 0$, a deterministic sign flip, while the remaining three ($0.606$, $0.910$, $1.668$) preserve the sign -- so a memoryless slicer errs on exactly one pattern in four as noise vanishes, giving the analytic floor of exactly $1/4$. The non-monotonicity follows from the same geometry: at moderate SNR, noise occasionally pushes the flipped sample back across the decision boundary, so the conditional error rate on that pattern rises toward one as SNR grows and total BER climbs toward $0.25$ from below -- standard behaviour for a closed-eye channel under symbol-by-symbol detection \cite{LuckySalzWeldon1968, ProakisSalehi2008DigitalComm}.

\subsection{\texorpdfstring{Complexity: the $M^{L}$ trellis against a fixed forward pass}{Complexity: the M to the L trellis against a fixed forward pass}}
MLSE maintains $M^{L-1}$ states and evaluates $M^L$ branch metrics for channel length $L$ \cite{Forney1972MLSE,ProakisSalehi2008DigitalComm}. The real $11\to13\to1$ MLP has $170$ parameters and $11\cdot13+13=156$ MACs, plus 14 bias additions and nonlinear activations. Counting a multiply and add as two FLOPs gives approximately 330 FLOPs excluding activation cost. These are not 312 MACs. Cost is fixed only for that deployed architecture; longer windows increase it. Reduced-state sequence estimation \cite{EyuboqluQureshi1988RSSE} is an unmeasured lower-cost classical alternative.

\subsection{The blind regime: CMA converges where decision-directed MLSE does not}

CMA is a blind equalizer for constant-modulus signals \cite{Godard1980SelfRecovering,TreichlerAgee1983CMA}. The measured random composite family also contains nonlinearity and phase uncertainty, so the zeros of one fixed linear tap vector do not establish invertibility or convergence for the whole cascade. The observed CMA and decision-directed MLSE results are implementation-specific. Documented convergence and decision-feedback issues \cite{Ding1991IllConvergence,RaheliPolydorosTzou1995PSP} motivate possible explanations but do not identify the cause of these measurements.

\subsection{The pilot-budget crossover: reliable at ten pilots, collapse at five}

A three-tap least-squares fit can be identifiable from five pilots if its design matrix has full column rank. For a correctly specified linear model with white noise, its covariance is $\sigma^2(\mathbf X^H\mathbf X)^{-1}$ \cite{Kay1993Estimation}; the simplification $\sigma^2 L/N_p$ additionally requires suitably normalized orthogonal pilot columns. The nonlinear composite channel need not satisfy that model. The observed small-pilot degradation is therefore not a measured rank threshold or an unbounded-error theorem. Other estimators, including meta-learned demodulation \cite{ParkSimeone2021MetaDemod}, are not benchmarked here.

\subsection{The finite-window gap is not causally decomposed}

For bit-MAP decisions under the same model and side information, giving the detector a full sequence cannot increase its minimum expected bit error relative to a window. This observation-ordering argument is not a bound on a particular Viterbi-versus-MLP BER gap, because Viterbi optimizes sequence error. No windowed bit-MAP ablation separates truncation, approximation, finite training, and optimization effects here. The corrected BPSK required-SNR penalties range from $0.26$ to $2.53$~dB at the stated targets. They cannot be attributed solely to the 170-parameter budget.

\paragraph{Supplementary numerical checks.} The across-seed spread in SNR penalty against DF in the separate native-architecture replication is $0.004$~dB for Mamba-S6, $0.007$ for the Hybrid, $0.013$ for Mamba-2, $0.043$ for the MLP, $0.075$ for the Transformer and $0.091$ for the VAE. These finite three-seed ranges are not confidence bounds on failure probabilities.

For the QPSK linear-MMSE detail run, the attained MSE at 16~dB is $0.0666,\allowbreak\ 0.0374,\allowbreak\ 0.0354,\allowbreak\ 0.0333$ for $N=3,5,7,11$. The required-SNR penalties at targets $10^{-1}$ and $10^{-2}$ are respectively $+3.24, +2.98, +2.79, +2.65$ and $+6.34, +4.27, +4.13, +4.02$~dB. These measurements do not imply monotonic BER at every target.
